\documentclass[10pt,a4paper]{amsart}
\usepackage[margin=19mm]{geometry}
\usepackage{amsmath,amssymb,mathtools}
\usepackage[scr=rsfs,cal=euler]{mathalfa}
\usepackage{xfrac}
\usepackage[unicode,hidelinks,bookmarks=false]{hyperref}
\newcommand{\ie}{i.e.\ }
\newcommand{\cf}{cf.\ }
\newcommand{\resp}{respectively\ }
\newcommand{\Subsection}{\subsection}
\providecommand{\nobreakdash}{\nobreak\hbox{-}}
\setlength{\emergencystretch}{2em}
\multlinegap0pt
\usepackage{amssymb}
\newtheorem{theorem}{Theorem}
\newtheorem{lemma}{Lemma}
\title{Prawitz's area theorem and \\ mixed Aharonov sequence}
\author{Jianjun Jin}
\date{Source: arXiv:2308.08081v4 (CC BY 4.0)}
\begin{document}
\maketitle

\noindent\emph{Source and licence.} Prawitz's area theorem and \\ mixed Aharonov sequence. Version arXiv:2308.08081v4 (CC BY 4.0), distributed by arXiv under the Creative Commons Attribution 4.0 International licence, http://creativecommons.org/licenses/by/4.0/, as recorded in the arXiv licence metadata for this exact version and shown on the abstract page https://arxiv.org/abs/2308.08081v4. Full licence notice: \texttt{licenses/arxiv-2308.08081-CC-BY-4.0.txt}.\par\smallskip

\noindent\textit{Editorial scope.} Counted unit: the article's theorem main-1 (source0.tex lines 313--318). The article's complete original proof of it is delivered with it (source0.tex lines 419--471, one \texttt{proof} environment of 53 lines, which the article places away from the statement). No mathematics is written, repaired or re-derived: the statement and the proof are the article's own lines, in the article's own notation, and no retained line is substituted or re-flowed.  Retained uncounted as the source-local context the proof needs: lines 171--171; lines 278--280; lines 327--331; lines 329--330; lines 346--352; lines 361--369.  Nothing was omitted from any retained span, so the package carries no omission record.  No retained line contains a citation, so the wrapper carries no bibliography and no bibliography entry.  No retained line contains a figure environment, an image-inclusion command or drawing code, so no image is delivered and the package declares no asset dependency.  Editorial formatting only, in addition: an AMS \texttt{amsart} preamble that restates the article's own declarations and macros used by the retained lines, namely the environments \texttt{theorem}, \texttt{lemma} on separate counters, together with the standard packages those lines need.  The retained environments print in this document as Theorem 1, Lemma 1 in order of appearance, whereas the article prints the same items with the numbering of its own full document.  The complete native source of the article as distributed in the arXiv e-print for this exact version is bundled unchanged as \texttt{sources/145694/source0.tex}.
\begin{equation}\label{sigma}\sigma_{z}(\zeta):=\frac{\zeta+z}{1+\bar{z}\zeta}, \, \zeta \in \mathbb{D}.\end{equation}

\begin{equation}\label{new}
\left[\frac{f'(z)(w-z)}{f(w)-f(z)}\right]^{\lambda}=\sum_{n=0}^{\infty}\Phi_{\lambda, n}(f; z)(w-z)^n.
\end{equation}

\begin{theorem}\label{main-3}
Let $\lambda>0$ and $f \in \mathcal{L}(\mathbb{D})$. Then $f$ is univalent analytic in $\mathbb{D}$ if and only if, for all $\zeta\in \mathbb{D}$, there is a positive constant $C(\lambda, \zeta)$ such that for all $n\geq 1$,
\begin{equation}\label{m-ine-3}\left|A_n(f; \zeta)\right|\leq C(\lambda, \zeta).
\end{equation}
\end{theorem}

\begin{equation}\left|A_n(f; \zeta)\right|\leq C(\lambda, \zeta).
\end{equation}

\begin{lemma}\label{lemma-1}
Let  $\lambda>0$ and $f\in \mathcal{U}(\mathbb{D})$. Then we have
 \begin{equation}\label{p-area}
 \sum_{n=1}^{\infty}(n-\lambda)|\Phi_n(f; 0)|^2\leq \lambda.
 \end{equation}
 Equality holds in (\ref{p-area}) if and only if $f$ is a full mapping.
\end{lemma}

\begin{lemma}\label{lemma-2}
Let $\lambda>0$ and $f\in \mathcal{L}(\mathbb{D})$. Let $w, \zeta\in \mathbb{D}$ be fixed and let $z=\sigma_{\zeta}(w)=\frac{w+\zeta}{1+\bar{\zeta}w}$.  We set
$$\mathbf{F}(w)=f(z).$$
Then we have 
\begin{eqnarray}\label{l-22}
\lefteqn{\Phi_{n}(\mathbf{F}; w)=} \nonumber\\
&&\sum_{j=0}^{n}(-1)^{n-j}\binom{\lambda}{n-j}\sum_{k=0}^{j}\binom{j-1}{j-k}\frac{(-\bar{\zeta})^{n-k}(1-|\zeta|^2)^{k}}{(1+\bar{\zeta}w)^{n+k}}\Phi_{k}(f; z), \, n\geq 0.
\end{eqnarray}
\end{lemma} 

\begin{theorem}\label{main-1}
Let $\lambda>0$ and $f \in \mathcal{L}(\mathbb{D})$. Then $f$ is univalent analytic in $\mathbb{D}$ if and only if, for all $\zeta\in \mathbb{D}$, it holds that
\begin{equation}\label{m-ine-1}
\sum_{n=1}^{\infty}(n-\lambda)\left|A_n(f; \zeta)\right|^2\leq \lambda.
\end{equation}
\end{theorem}

\begin{proof}[Proof of Theorem \ref{main-1} and \ref{main-3}]
We first prove the only if part. We assume that $f\in \mathcal{U}(\mathbb{D})$.  Let $w, \zeta\in \mathbb{D}$ be fixed and let $z=\frac{w+\zeta}{1+\bar{\zeta}w}$.  We set $\mathbf{F}(w)=f(z)$.  Then we have 
\begin{eqnarray}\label{pro-1}
\lefteqn{\Phi_{n}(\mathbf{F}; 0)=} \\
&&\sum_{j=0}^{n}(-1)^{n-j}\binom{\lambda}{n-j}\sum_{k=0}^{j}\binom{j-1}{j-k}(-\bar{\zeta})^{n-k}(1-|\zeta|^2)^{k}\Phi_{k}(f; \zeta),\, n \geq 0, \nonumber
\end{eqnarray}
by Lemma \ref{lemma-2}, and \begin{equation}\label{pro-2}
 \sum_{n=1}^{\infty}(n-\lambda)|\Phi_n(\mathbf{F}; 0)|^2\leq \lambda,
 \end{equation}
 by  Lemma \ref{lemma-1} since $\mathbf{F}$ is also univalent analytic in $\mathbb{D}$. Thus (\ref{m-ine-1}) follows by (\ref{pro-1}) and (\ref{pro-2}) and the only if part of Theorem \ref{main-1} is proved.
 
 When $\lambda\in (0,1)$, (\ref{pro-2}) implies that  
\begin{equation}|\Phi_n(\mathbf{F}; 0)|\leq  \sqrt{\frac{\lambda}{1-\lambda}},  \nonumber
\end{equation}for all $\zeta\in \mathbb{D}, n\geq 1.$

When $\lambda\geq1$,  there is a positive integer $\mathbf{N}$ such that $\mathbf{N}\leq \lambda < \mathbf{N}+1$. Then (\ref{pro-2}) yields  that  
\begin{eqnarray}
|\Phi_n(\mathbf{F}; 0)| \leq\sqrt{\frac{\lambda+ \mathbf{E}(\lambda, \zeta)}{\mathbf{N}+1-\lambda}}:=\mathbf{E}_1(\lambda, \zeta), \nonumber
\end{eqnarray}
for all $\zeta\in \mathbb{D}, n\geq \mathbf{N}+1$.  Here $$ \mathbf{E}(\lambda, \zeta)=\sum_{k=1}^{N}(\lambda-k)|\Phi_k(\mathbf{F}; 0)|^2.$$
Meanwhile, it's obvious that 
$$|\Phi_n(\mathbf{F}; 0)|  \leq \sqrt{\sum_{k=1}^{\mathbf{N}}|\Phi_k(\mathbf{F}; 0)|^2}:=\mathbf{E}_2(\lambda, \zeta),$$ for all $n \in [1, \mathbf{N}].$ 

Then, from (\ref{pro-1}), we find that (\ref{m-ine-3}) is true for all $\zeta\in \mathbb{D}, n\geq 1$,  when we take 
$$C(\lambda, \zeta)=\sqrt{\frac{\lambda}{1-\lambda}},$$ for $\lambda\in (0,1)$, 
and  
$$C(\lambda, \zeta)=\mathbf{E}_1(\lambda, \zeta)+\mathbf{E}_2(\lambda, \zeta),$$
for $\lambda\geq 1$.  The only if part of Theorem \ref{main-3} is proved.

We proceed to prove the  if part.  We assume that (\ref{m-ine-3}) holds. If $f$ is not univalent analytic in $\mathbb{D}$, then there are two different points $\zeta, \alpha$ in $\mathbb{D}$ with $f(\zeta)=f(\alpha)$. We let
$$\mathcal{F}(z)=f(\sigma_{\zeta}(z))\,\, {\text{and}}\,\ \alpha=\sigma_{\zeta}(\beta).$$
Here $\sigma_{\zeta}$ is defined as in (\ref{sigma}). Then we get that 
\begin{equation}\label{proof}\mathcal{F}(0)=f(\zeta)=f(\alpha)=\mathcal{F}(\beta). \end{equation}
It is easy to see that $\beta \neq 0$ since $\zeta \neq \alpha$.

On the other hand, from (\ref{new}), we have 
\begin{equation}\label{proof-1}
\left[\frac{{\mathcal{F}}'(0)\beta}{\mathcal{F}(\beta)-\mathcal{F}(0)}\right]^{\lambda}=1+\sum_{n=1}^{\infty}\Phi_n(\mathcal{F}; 0){\beta}^n.
\end{equation}
Also, we see from the assumption and (\ref{pro-1}) that
\begin{equation}\label{proof-2}
|\Phi_{n}(\mathcal{F}; 0)|\leq C(\lambda, \zeta),
\end{equation} 
for all $n\geq 1.$
It follows from (\ref{proof-1}) and (\ref{proof-2}) that  
\begin{equation}
\frac{|{\mathcal{F}}'(0)\beta|^{\lambda}}{|\mathcal{F}(\beta)-\mathcal{F}(0)|^{\lambda}}\leq 1+\sum_{n=1}^{\infty}|\Phi_n(\mathcal{F}; 0)|{|\beta|}^n\leq 1+C(\lambda, \zeta)\frac{|\beta|}{1-|\beta|} <+\infty. \nonumber
\end{equation}
This contradicts (\ref{proof}) so that $f$ is univalent analytic in $\mathbb{D}$.  The if part of Theorem \ref{main-3} is proved and the proof of Theorem \ref{main-3} is finished.

Next, we assume that (\ref{m-ine-1}) holds, then we know from the above arguments that  (\ref{m-ine-3}) holds with $C(\lambda, \zeta)=\sqrt{\lambda/(1-\lambda)}$ when $\lambda\in (0,1)$, 
and  $C(\lambda, \zeta)=\mathbf{E}_1(\lambda, \zeta)+\mathbf{E}_2(\lambda, \zeta)$ when $\lambda\geq 1.$  By using the if part of Theorem \ref{main-3},  we see that $f$ is univalent analytic in $\mathbb{D}$. The proof of Theorem \ref{main-1} is complete. 
\end{proof}

\end{document}
